% ===================================================================
%  పట్టిక సంఖ్య 7 — చలికాలంలో పల్లె. వసంతంలో పల్లెటిల్లు.
%
%  Traduction, faite en 2026, en telougou d'aujourd'hui, sur l'ido de
%  texto/io. Regles de la colonne : voir 10-pattika-01.tex. Deux
%  scenes, deux numerotations : les cles portent c1 et c2. Le
%  feuillet 49 porte aussi la ligne d'apparat « రెండవ శ్రేణి ».
%
%  LE GARDE CHAMPETRE S'APPELLE తలారి, ET C'EST LA TROISIEME FOIS
%  QUE CE POSTE TROUVE SON NOM. Le marathi disait कोतवाल, l'officier
%  du village marathe ; le telougou dit తలారి, le veilleur de nuit du
%  village en pays andhra — une charge de village reelle, en fonction
%  bien avant le gendarme du livret. Deux langues sans parente, deux
%  administrations villageoises differentes, et la meme fonction
%  nommee sans periphrase. C'est le pendant exact du బంట్రోతు du
%  tableau 1, qui repondait au فرّاش egyptien et au शिपाई marathe.
%
%  LE VIEUX CORDONNIER FAIT LA MEME BLAGUE, ET C'EST LA QUATRIEME
%  COLONNE. « En ville je m'appelais shuifisto et je n'avais pas de
%  clients ; ici je m'appelle shureparisto et le travail ne manque
%  pas. » L'egyptien opposait le savant صانع أحذية au turc جزمجي,
%  le marathi le sanskrit पादत्राणनिर्माता au nom de caste चांभार. Le
%  telougou oppose పాదరక్షల తయారీదారు a చెప్పులు కుట్టేవాడు — le
%  fabricant de chaussures contre celui qui les coud.
%
%  ET L'ON N'ECRIT PAS LE NOM DE CASTE ICI. Le marathi pouvait dire
%  चांभार, qui est devenu un nom de famille ordinaire ; les noms
%  telougous correspondants designent des castes opprimees et ne
%  s'emploient pas en 2026 comme etiquettes de metier. On ecrit donc
%  la fonction — చెప్పులు కుట్టేవాడు, గొర్రెల కాపరి —, et la blague
%  du cordonnier tient quand meme, parce qu'elle repose sur le
%  registre et non sur la caste. Ecrire en telougou d'aujourd'hui,
%  c'est aussi cela.
%
%  LA FERME EST ENTIEREMENT TELOUGOUE, ET C'EST LE PAYS DU RIZ QUI
%  PARLE :
%
%    పొలం (3)     la campagne     చేను (4)     le champ
%    నాగలి (5)    la charrue      చాలు (8)     le sillon
%    గొర్రు (40)  la herse        గునపం (39)   la pioche
%    కొడవలి       la faucille     కంకులు       les epis
%    కళ్ళం (33)   l'aire          బావి (29)    le puits
%    కొట్టం (58)  l'etable        మంద (49)     le troupeau
%    పాలేరు (65)  le valet de ferme
%    రాట్నం       le rouet        కదురు        le fuseau
%    కవ్వం (61)   la baratte      పెంట (81)    le fumier
%    ఆబోతు (27)   le taureau      దూడ (24)     le veau
%    సిగ (72)     la crete du coq నాడా (6)     le fer a cheval
%    డెక్క        le sabot        వాగు (23)    le ruisseau
%    సత్రం (2)    l'auberge       శాలువా (43)  le chale
%    చుట్ట        la pipe — le cheroot, qui est ce qu'on fume ici
%
%  ET LE TROUPEAU SE NOMME BETE PAR BETE, comme en marathe :
%  పొట్టేలు le belier, గొర్రె la brebis, గొర్రెపిల్ల l'agneau, మేక la
%  chevre, మేకపిల్ల le chevreau.
%
%  UN EMPRUNT ASSUME, LE MEME QU'EN MARATHE : le marron chaud. Le
%  chataignier ne pousse pas plus en Andhra qu'au Maharashtra ; on
%  ecrit చెస్ట్‌నట్లు et l'on ne transpose pas la chose.
%
%  ECRIT AVEC LA LISTE DE parigi.py EN MAIN : trente-et-une paires
%  signalees, dix-sept vraies, dont deux tenues par un PARTICIPE
%  — « le soleil inondant la campagne de ses rayons », « le foin
%  seche en le retournant a la fourche » —, que la premiere version
%  du motif aurait manquees.
% ===================================================================

%%K t07-noto-1 noto
\VUnotes{143.2mm}{%
(*) భోజన వివరాల కోసం పట్టికలు 6 మరియు 13 చూడండి.%
}

%%K t07-apar-1 apar
\VUsaut{4.0mm}
\VUcentre{13.2pt}{90}{రెండవ శ్రేణి}
\VUsaut{3.0mm}
\VUfilet{16mm}
\VUsaut{5.6mm}
\VUtitre{15.0pt}{60}{పట్టిక సంఖ్య 7}
\VUsaut{3.0mm}

%%K t07-tit-1 sub
\VUcentre{12.0pt}{0}{\textit{మొదటి రంగం.}}
\VUsaut{3.0mm}

%%K t07-tit-2 sub
\VUpk{10.2pt}{40}{చలికాలంలో పల్లె.}
\VUsaut{4.0mm}

%%K t07-c1-01-1 p
1. --- నూతన సంవత్సర సెలవుల్లో నేను మా నాన్నతో నవనగరం వెళ్ళాను — అది ఒక చిన్న ఊరు, అక్కడ ఆయనకు
కొన్ని వ్యాపార పనులు ఉన్నాయి.

%%K t07-c1-02-1 p
2. --- పట్టణానికీ ఆ ఊరికీ మధ్య నడిచే \VUgras{గుర్రపు బండి} \textsuperscript{(11)} మేము ఎక్కాం;
కానీ చాలా చలిగా ఉండటంతో అంత సౌకర్యం లేని ఆ బండిలో ప్రయాణం పెద్ద సుఖంగా అనిపించలేదు.

%%K t07-c1-03-1 p
3. --- అందుకే \VUgras{బండివాడు} \textit{తెల్ల ఎలుగుబంటి} \VUgras{సత్రం} \textsuperscript{(2)}
ముందున్న చౌకలో బండి ఆపడం చూసి మాకు నిజమైన ఆనందం కలిగింది. \VUgras{సత్రం యజమాని}
\textsuperscript{(4)} తన ఇంటి గడపపై నిమ్మళంగా \VUgras{చుట్ట} కాలుస్తూ మా కోసం
ఎదురుచూస్తున్నాడు.

%%K t07-c1-03-2 p
ఆయన మమ్మల్ని లోపలికి పిలిచారు, పొయ్యిలో చిటపటలాడుతున్న ఎండుకొమ్మల మంట ముందు కూర్చునేందుకు నేను
ఏమాత్రం మొహమాటపడలేదు. అప్పుడు పాత \VUgras{బోర్డును} \textsuperscript{(3)} కిర్రుమనిపిస్తూ,
పొగగొట్టం నుంచి వచ్చే \VUgras{పొగను} \textsuperscript{(1)} నల్లని సుడులతో ఎగరగొడుతున్న ఆ కోసే
\VUgras{ఉత్తరపు గాలిని} నేను బాగానే ఎగతాళి చేశాను.

%%K t07-c1-04-1 p
4. --- భోజన సమయమైంది. ఎక్కువ సేపు ఎదురుచూడకుండా, మాకు వడ్డించిన సాదా కానీ రుచికరమైన భోజనం మేము
తృప్తిగా తిన్నాం (*).

%%K t07-tit-3 sub
\VUpk{10.2pt}{40}{కొన్ని సందర్శనలు.}
\VUsaut{3.0mm}

%%K t07-c1-05-1 p
5. --- భోజనం తర్వాత నేను మా నాన్నతో వెళ్ళాను — ఆయన బీమా చేసేవారు — ఆయన ఖాతాదారుల దగ్గరికి.
ముందుగా మేము \VUgras{కమ్మరి} \textsuperscript{(5)} దగ్గరికి వెళ్ళాం, అతను సత్రం దగ్గరే ఉంటాడు.
అతను గుర్రానికి నాడా కొట్టడంలో మునిగి ఉన్నాడు. అతని సహాయకుడి బలం చూసి నాకు ఆశ్చర్యం వేసింది —
గుర్రపు \VUgras{డెక్కను} అతను తన తోలు అప్రాన్‌పై గట్టిగా పట్టుకున్నాడు, ఈలోగా కమ్మరి బలమైన
దెబ్బలతో \VUgras{నాడా} \textsuperscript{(6)} బిగించాడు.

%%K t07-c1-06-1 p
6. --- ఆ తర్వాత మేము ఊరి \VUgras{చెప్పులు కుట్టేవాడి} \textsuperscript{(7)} దగ్గరికి వెళ్ళాం —
ముసలి జేకబ్. బోర్డుగా అతని దగ్గర ఒక అందమైన \VUgras{బూటు} ఉన్నా, చిల్లులు పడిన పాత చెప్పులకు
కొత్త అరికాళ్ళు వేయడంతోనే అతను తృప్తిపడేవాడు.

%%K t07-c1-06-2 p
ఆ ముసలాయన నవ్వుతూ మాతో అన్నాడు : « పట్టణంలో నేను « పాదరక్షల తయారీదారు »ని, నాకు ఖాతాదారులు
లేరు. ఇక్కడ నేను « చెప్పులు కుట్టేవాడి »ని, పనికి కొదవ లేదు. »

%%K t07-c1-07-1 p
7. --- తన పొరుగున ఉన్న \VUgras{మంగలి} \textsuperscript{(8)} గురించి అతను మాకు చెప్పాడు — అతని
ఖాతాదారులు అసంతృప్తిగా ఉంటారు. అతని \VUgras{కత్తులు} ఎప్పుడూ చిల్లులు పడి ఉంటాయి, అతని
\VUgras{కత్తెరలు} ఎప్పుడూ సరిగా కత్తిరించవు. కానీ అతనే ఊరిలో ఒక్కగానొక్క మంగలి కావడంతో, గడ్డమూ
జుట్టూ అతని దగ్గరే చేయించుకోక తప్పదు, అది స్పష్టం. పైగా అతను పిసినారి, ముభావపు మనిషి.

%%K t07-c1-08-1 p
8. --- అదృష్టవశాత్తూ మిగతా \VUgras{పొరుగువాళ్ళు} అతనిలా లేరు. ఉదాహరణకు
\VUgras{బండ్లు చేసేవాడు} \textsuperscript{(9)} మంచి మనిషి, కష్టపడేవాడు, సాయపడేవాడు. అతని దగ్గర
బండి బాగుచేయించడం ఆనందం. అతని పని ఎప్పుడూ చక్కగా పూర్తయి ఉంటుంది.

%%K t07-c1-09-1 p
9. --- ముసలి జేకబ్ \VUgras{జీను చేసేవాడి} \textsuperscript{(10)} గురించి కూడా ఫిర్యాదు
చేయలేదు; పని ఒత్తిడిగా ఉన్నప్పుడు అతను అప్పుడప్పుడూ \VUgras{గుర్రపు సామాను} కుట్టడంలో సాయం
చేస్తాడు.

%%K t07-c1-09-2 p
మా నాన్న అతని ఇంటివాళ్ళ గురించి అడిగారు : « అందరూ బాగున్నారు. నా మనవరాలు \VUgras{బడిలో}
\textsuperscript{(12)} ఉంది. ఆమె \VUgras{ఉపాధ్యాయిని} \textsuperscript{(13)} ఆమెపై చాలా
సంతోషంగా ఉంది, ఆమె మంచి \VUgras{విద్యార్థిని} \textsuperscript{(14)}. నా చెడిపోయిన
కొడుకులాంటిది కాదు; వాడికి ఆటలే ధ్యాస. \VUgras{ఉపాధ్యాయుడు} \textsuperscript{(15)} వాడికి ఏదీ
నేర్పలేకపోతున్నాడు. »

%%K t07-tit-4 sub
\VUsaut{3.0mm}
\VUpk{10.2pt}{40}{పల్లె కబుర్లు.}
\VUsaut{3.0mm}

%%K t07-c1-10-1 p
10. --- ఆ తర్వాత ఆ ముసలాయన మాకు ఊరి కబుర్లు చెప్పాడు. కొత్త బడి కట్టబోతున్నారు, ఎందుకంటే
\VUgras{పంచాయతీ భవనంలోని} బడి చాలా చిన్నదైపోయింది. పైగా అది సులువే, ఎందుకంటే
\VUgras{మిల్లువాడి} తోటలో ఒక ముక్క కొంటే సరిపోతుంది. ఆ పాపం మనిషికి దానివల్ల చాలా లాభం.
చలివల్ల \VUgras{నీటిమిల్లు} \textsuperscript{(16)} \VUgras{తిరగలి రాళ్ళు} గోధుమలు
విసరలేకపోతున్నాయి, ఎందుకంటే \VUgras{చక్రం} \textsuperscript{(17)} \VUgras{గడ్డకట్టిన మంచుతో}
\textsuperscript{(25)} ఆగిపోయింది — \VUgras{వాగులోని} \textsuperscript{(23)} మంచుతో.

%%K t07-c1-11-1 p
11. --- « కట్టడాల సంగతంటే, మా కొత్త \VUgras{చర్చి} \textsuperscript{(18)} చూశారా? మంచు దుప్పటి
కింద అది చాలా చిత్రంగా ఉంది. పాత గంటస్తంభాన్ని ఇక గుర్తుపట్టలేక \VUgras{కాకులు}
\textsuperscript{(19)} \VUgras{సమాధి స్థలంలోని} \textsuperscript{(20)} పెద్ద చెట్టుపై దిగాలుగా
వాలతాయి. »

%%K t07-c1-12-1 p
12. --- సమాధి స్థలం ప్రస్తావనతో చెప్పులు కుట్టేవాడికి మా నాన్నకు తెలిసిన, గత సంవత్సరం చనిపోయిన
మనుషులు గుర్తొచ్చారు. « ఎన్ని \VUgras{సమాధులు} \textsuperscript{(21)}, అయ్యా,
\VUgras{సైప్రస్ కింద} \textsuperscript{(22)} మిగతావాటికి తోడయ్యాయి! మృత్యువు మా ముసలి నోటరీని
కోసుకుపోయింది. ఊరంతా ఆయన \VUgras{అంత్యక్రియలకు} వచ్చింది. »

%%K t07-c1-13-1 p
13. --- « అదిగో ఆయన వారసుడు, వాగుపై జారుతున్నాడు. తన \VUgras{స్కేటు} \textsuperscript{(26)}
తెగిన పట్టీ నా దగ్గర బాగుచేయించుకునేందుకు ఇప్పుడే వచ్చి వెళ్ళాడు. »

%%K t07-c1-14-1 p
14. --- « ఇదిగో వాగు ఒడ్డున కొత్త \VUgras{తలారి} \textsuperscript{(27)}. అతను నివృత్త పోలీసు,
ఊరి కొంటె పిల్లలకు భయం పుట్టిస్తాడు. అతని \VUgras{కాళ్ళ పట్టీలూ} \textsuperscript{(29)} అతని
\VUgras{టోపీ} \textsuperscript{(28)} కనిపించగానే వాళ్ళు వెంటనే పరుగు లంకించుకుంటారు. »

%%K t07-c1-15-1 p
15. --- « అరే! ఇదిగో ముసలి యోసేపు, రొట్టెలవాడికి \VUgras{కట్టెల బండి} \textsuperscript{(30)}
తోలుకెళ్తున్నాడు. --- నిజమే, అన్నారు మా నాన్న, అతని \VUgras{కంచరగాడిదల} \textsuperscript{(31)}
జోడి నాకు తెలుసు. అతనితో నాకు మాట్లాడాల్సిందే ఉంది, ఈ అవకాశం వాడుకుంటాను. వస్తాను, జేకబ్ తాతా.
»

%%K t07-c1-16-1 p
16. --- మేము బయటికి వచ్చేసరికి ఊరి పిల్లలు బడి వదిలి \VUgras{జారుడు బండ}
\textsuperscript{(33)} వైపు పరుగెత్తారు — \VUgras{పడి} \textsuperscript{(34)} చేయి కాలు
విరగ్గొట్టుకునే ప్రమాదం లెక్కచేయకుండా. వాళ్ళలో ఒకడు బండ్లు చేసేవాడి దగ్గర తన
\VUgras{జారుడు బండి} \textsuperscript{(32)} తీసుకుని, దానిలో తన సహచరుడిని కూర్చోబెట్టి
పరుగెత్తాడు.

%%K t07-c1-17-1 p
17. --- మిగతావాళ్ళు \VUgras{మంచు కుప్పవైపు} \textsuperscript{(37)} పరుగెత్తారు —
\VUgras{వంతెన} \textsuperscript{(40)} దగ్గరిది — మరియు నిన్ననే నిలబెట్టిన
\VUgras{మంచు బొమ్మపై} \textsuperscript{(39)} దాడి మొదలుపెట్టారు; దానికి వాళ్ళు టోపీ, చుట్ట,
భుజాన బడి సంచి తగిలించారు.

%%K t07-c1-18-1 p
18. --- గడ్డకట్టించే గాలినీ చలినీ వాళ్ళు లెక్కచేయరు. చాలామంది మెడలో \VUgras{మఫ్లర్}
\textsuperscript{(35)} కూడా లేదు, \VUgras{మంచు ఉండల} \textsuperscript{(38)} యుద్ధం తర్వాత
మళ్ళీ వెచ్చబడేందుకు వాళ్ళకు గుప్పెడు వేడివేడి \VUgras{చెస్ట్‌నట్లు} \textsuperscript{(42)}
చాలు — ముసలి \VUgras{అమ్మేది} జాకోబినా దగ్గర కొన్నవి; ఆమె తన చాలా పలచని \VUgras{శాలువా}
\textsuperscript{(43)} కింద చలికి వణుకుతోంది.

%%K t07-tit-5 sub
\VUsaut{3.0mm}
\VUcentre{12.0pt}{0}{\textit{రెండవ రంగం.}}
\VUsaut{3.0mm}

%%K t07-tit-6 sub
\VUpk{10.2pt}{40}{వసంతంలో పల్లెటిల్లు.}
\VUsaut{4.0mm}

%%K t07-c2-01-1 p
1. --- చలికాలపు మజా అంతా ఉన్నా, \VUgras{వసంతం} తిరిగి రావడాన్ని మనం మనస్ఫూర్తిగా ఆనందంతో
స్వాగతిస్తాం.

%%K t07-c2-01-2 p
మే నెలలో, సాయంత్రం, పొలం అంగణం ఎంత ఆనందకరమైన చిత్రంలా కనిపిస్తుంది!

%%K t07-c2-02-1 p
2. --- క్షితిజంపై \VUgras{సూర్యుడు} \textsuperscript{(1)} నెమ్మదిగా అస్తమిస్తున్నాడు,
\VUgras{పొలాన్ని} \textsuperscript{(3)} తన ఊదారంగు \VUgras{కిరణాలతో} \textsuperscript{(2)}
ముంచెత్తుతూ. కష్టపడే \VUgras{దున్నేవాడు} \textsuperscript{(6)} తన \VUgras{చేను}
\textsuperscript{(4)} దున్నేందుకు తొందరపడుతున్నాడు.

%%K t07-c2-02-2 p
తన \VUgras{నాగలితో} \textsuperscript{(5)} — దాన్ని బలమైన \VUgras{గుర్రానికి}
\textsuperscript{(7)} కట్టారు — అతను \VUgras{చాలు} \textsuperscript{(8)} తీస్తాడు; ఆ చాలులో
\VUgras{విత్తనం చల్లేవాడు} \textsuperscript{(9)} గుప్పెడు \VUgras{విత్తనాలు} వేస్తాడు,
వాటినుంచి బంగారు \VUgras{కంకులు} వస్తాయి.

%%K t07-c2-03-1 p
3. --- \VUgras{గడ్డి కోసేవాళ్ళు} \textsuperscript{(11)} తమ కొడవళ్ళతో బీడులోని గడ్డి కోశారు,
ఎండుగడ్డి చేసేవాళ్ళు \VUgras{ఎండుగడ్డిని} \textsuperscript{(10)} ఆరబెట్టారు —
\VUgras{గడ్డి పారలతో} \textsuperscript{(12)} మరియు దంతెలతో తిరగేస్తూ —, ఇదిగో వాళ్ళు తిరిగి
వస్తున్నారు.

%%K t07-c2-03-2 p
కొందరు ఎద్దులు లాగే బరువైన బండి ముందు నడుస్తున్నారు; తమ పనిముట్టు భుజాన మోసుకెళ్తున్నారు.
ఇంకొందరు, కాస్త బద్ధకస్తులు, సువాసనగల ఎండుగడ్డిపై హాయిగా పడుకున్నారు.

%%K t07-c2-04-1 p
4. --- వెళ్తూవెళ్తూ వాళ్ళు \VUgras{తోటమాలికి} \textsuperscript{(16)} స్నేహంగా చేయి ఊపుతారు —
అతను \VUgras{పండ్ల తోటలో} \textsuperscript{(13)} \VUgras{పండ్ల చెట్లు} కత్తిరించడంలో,
\VUgras{బేరి పండ్ల చెట్లకు} \textsuperscript{(14)} అంటుకట్టడంలో మునిగి ఉన్నాడు.
\VUgras{ప్లం చెట్లు} \textsuperscript{(15)} పూత పట్టాయి, బాగా ఎరువు వేసిన
\VUgras{కూరగాయల తోటలో} \textsuperscript{(17)} కూరగాయలూ \VUgras{క్యాబేజీలూ} కోసంబరి ఆకులూ
ఇప్పటికే బాగా పెరిగాయి.

%%K t07-c2-04-2 p
గోడ కంచెపై ఒక అందమైన \VUgras{నెమలి} \textsuperscript{(18)} తన పింఛం విప్పుతోంది, తన అందమైన
ఈకలు చూపేందుకు.

%%K t07-tit-7 sub
\VUpk{10.2pt}{40}{పొలం అంగణం.}
\VUsaut{3.0mm}

%%K t07-c2-05-1 p
5. --- \VUgras{గుర్రపుశాల} \textsuperscript{(19)} తలుపులు తెరిచి ఉన్నాయి, లోపల మేత తొట్టి
మరియు \VUgras{పరుపుగడ్డి} \textsuperscript{(23)} కనిపిస్తాయి — \VUgras{ఆడగుర్రానివి}
\textsuperscript{(21)}, దాన్ని \VUgras{గుర్రపు కాపరి} \textsuperscript{(20)} ఇప్పుడే జోడు
నుంచి విప్పాడు. తన పొడవాటి కాళ్ళతో అంత తమాషాగా కనిపించే \VUgras{గుర్రపు పిల్ల}
\textsuperscript{(22)} గంతులు వేస్తూ తన తల్లి వెంట వెళ్తోంది.

%%K t07-c2-06-1 p
6. --- \VUgras{బావి} \textsuperscript{(29)} దగ్గర \VUgras{ఆవులు} \textsuperscript{(25)} నీళ్ళు
తాగేందుకు చెక్క \VUgras{తొట్టిలోకి} \textsuperscript{(26)} వెళ్తాయి, అది వాటికి నీటిచోటుగా
వాడతారు. \VUgras{దూడ} \textsuperscript{(24)} ఈ అవకాశం చూసుకుని పాలు తాగుతోంది, కొట్టం మంచి
వాసన పీలుస్తూ \VUgras{ఆబోతు} \textsuperscript{(27)} ఆనందంగా రంకె వేసి దృఢమైన \VUgras{కొమ్ముల}
\textsuperscript{(28)} తల గర్వంగా ఎత్తుతోంది.

%%K t07-c2-07-1 p
7. --- అందరూ పనిలో ఉన్నారు. \VUgras{పనిమనిషి} \textsuperscript{(32)} బావి \VUgras{తాడు}
\textsuperscript{(31)} చేత్తో పట్టుకుంది. పశువులకు నీళ్ళు పెట్టేందుకు ఆమె \VUgras{బకెట్‌తో}
\textsuperscript{(30)} నీళ్ళు తోడుతోంది. \VUgras{కళ్ళం} \textsuperscript{(33)} దగ్గర ఒక మనిషి
గడ్డి పోగుచేస్తున్నాడు, \VUgras{ఇంటి} \textsuperscript{(34)} తలుపు ముందు ఒక ముసలి
\VUgras{నూలు వడికేది} \textsuperscript{(35)} తన రాట్నంతో, తన \VUgras{కదురుతో} పాతకాలంలోలాగే
\VUgras{అవిసె} లేదా \VUgras{జనుము} వడుకుతోంది.

%%K t07-tit-8 sub
\VUsaut{3.0mm}
\VUpk{10.2pt}{40}{రైతు.}
\VUsaut{3.0mm}

%%K t07-c2-08-1 p
8. --- \VUgras{రైతు} \textsuperscript{(36)} ఈ పనులన్నీ నడిపిస్తాడు, తన పాలేర్లకు సూచనలు
ఇస్తాడు. \VUgras{గొర్రు} \textsuperscript{(40)} మరియు \VUgras{గునపం} \textsuperscript{(39)}
కప్పు కిందికి తీసుకోమని అతను చెబుతాడు — అవి \VUgras{గూడు} \textsuperscript{(38)} దగ్గర
మరిచిపోయారు, \VUgras{కుక్క} \textsuperscript{(37)} గూడు దగ్గర.

%%K t07-c2-09-1 p
9. --- అతని కేకలకు ఒక \VUgras{పావురం} \textsuperscript{(41)} భయపడుతుంది, అది రెక్కలు
టపటపలాడిస్తూ \VUgras{పావురాల గూట్లోకి} \textsuperscript{(42)} దూరుతుంది; \VUgras{కోళ్ళు}
\textsuperscript{(43)} కూడా, అవి \VUgras{కోళ్ళ గూట్లోకి} \textsuperscript{(44)} తిరిగి
వెళ్ళేందుకు తొందరపడతాయి.

%%K t07-c2-10-1 p
10. --- \VUgras{గొర్రెల దొడ్డి} \textsuperscript{(48)} ముందు, ఇదిగో ముసలి
\VUgras{గొర్రెల కాపరి} \textsuperscript{(45)}, తన \VUgras{మందను} \textsuperscript{(49)} తిరిగి
తోలుకొస్తున్నాడు. తన \VUgras{కర్ర} \textsuperscript{(46)} చేతిలో పట్టుకుని — అది అతని దగ్గర
లేకుండా ఉండదు, అలాగే అతని రంగు వెలసిన \VUgras{అంగీ} \textsuperscript{(47)} కూడా — అతను కొట్టం
వైపు \VUgras{పొట్టేళ్ళను} \textsuperscript{(50)}, \VUgras{గొర్రెలను} \textsuperscript{(53)},
\VUgras{గొర్రెపిల్లలను} \textsuperscript{(54)}, \VUgras{మేకలను} \textsuperscript{(51)} మరియు
\VUgras{మేకపిల్లలను} \textsuperscript{(52)} తోలుకెళ్తాడు. అతని విశ్వాసపాత్రమైన \VUgras{కుక్క}
\textsuperscript{(55)} ఆర్గుస్ అందరికంటే చివర వస్తుంది, మొరుగుతూ వెనుకబడినవాటిని
తొందరపెడుతుంది.

%%K t07-c2-11-1 p
11. --- ఈ గోలంతా, ఈ కదలికంతా ఆ మంచి \VUgras{పాడి ఆవు} \textsuperscript{(56)} నిమ్మళాన్ని
చెడగొట్టడం లేదు — అది తన \VUgras{గంట} \textsuperscript{(57)} గణగణలాడిస్తోంది, ఈలోగా
\VUgras{కొట్టం} \textsuperscript{(58)} తలుపు ముందు దాని పాలు పితుకుతున్నారు.

%%K t07-c2-12-1 p
12. తాను పాలు ఇవ్వాలని దానికి తెలిసినట్టుంది — అప్పుడే \VUgras{రైతమ్మ} \textsuperscript{(60)}
\VUgras{కవ్వంతో} \textsuperscript{(61)} \VUgras{వెన్న} తీయగలదు, లేదా ఆ చక్కని \VUgras{చీజ్}
\textsuperscript{(64)} చేయగలదు — అవి \VUgras{పెద్ద తొట్టిపై} \textsuperscript{(63)} పెట్టిన
అరపై నుంచి కారుతూ ఉంటాయి.

%%K t07-c2-13-1 p
13. --- \VUgras{పాలగదిని} \textsuperscript{(59)} అంతా \VUgras{పాలేరు} \textsuperscript{(65)}
చాలా శుభ్రంగా ఉంచుతాడు — అతను ఇప్పుడే \VUgras{పాల గిన్నెలు} \textsuperscript{(62)} చేత్తో
మోసుకొచ్చిన పెద్ద బకెట్‌లో కడిగాడు.

%%K t07-tit-9 sub
\VUsaut{3.0mm}
\VUpk{10.2pt}{40}{కోళ్ళ దొడ్డి.}
\VUsaut{3.0mm}

%%K t07-c2-14-1 p
14. --- తన మందపాటి చెక్క చెప్పులతో అతను కాస్త బరువుగా నడుస్తాడు, దాంతో \VUgras{టర్కీ కోడి}
\textsuperscript{(68)}, \VUgras{ఆడ బాతులు} \textsuperscript{(66)}, \VUgras{మగ బాతులు}
\textsuperscript{(67)} మరియు \VUgras{హంసలు} \textsuperscript{(69)} పారిపోతాయి — ఆ హంసలు
\VUgras{తమ ముక్కు} \textsuperscript{(70)} బార్లా తెరిచి కంగారుగా అరుస్తాయి.

%%K t07-c2-14-2 p
\VUgras{కోడిపుంజు} \textsuperscript{(71)}, కాస్త ఎక్కువ పోట్లాటకారి, తన ఎర్రని \VUgras{సిగ}
\textsuperscript{(72)} నిక్కబొడుచుకుంటుంది, తన \VUgras{తోక} \textsuperscript{(73)} ఈకలు
దులుపుకుంటుంది, « కొక్కొరోకో » అని గట్టిగా కూస్తుంది.

%%K t07-c2-15-1 p
15. --- \VUgras{తల్లికోడి} \textsuperscript{(74)} \VUgras{పెంటపై} \textsuperscript{(81)}
గాలిస్తోంది, తన \VUgras{కోడిపిల్లలతో} \textsuperscript{(75)}. \VUgras{పంది పిల్ల}
\textsuperscript{(78)} తన తల్లి దగ్గరికి పరుగెడుతోంది — ఆ అపారమైన \VUgras{ఆడ పంది}
\textsuperscript{(77)} దగ్గరికి, అది పందుల దొడ్డి దగ్గర మనకు కనిపిస్తోంది.

%%K t07-c2-16-1 p
16. --- తమ గదుల్లో \VUgras{కుందేళ్ళు} \textsuperscript{(79)} రైతు కూతురు తెచ్చిన లేత
\VUgras{గడ్డిని} \textsuperscript{(80)} ఆబగా తింటున్నాయి. జానినాకు తన కుందేళ్ళంటే చాలా ఇష్టం,
రోజూ కోళ్ళ గూట్లో తాజా \VUgras{గుడ్లు} \textsuperscript{(76)} తెచ్చాక ఆమె తన చిన్ని
స్నేహితులను చూసేందుకు వస్తుంది.

\VUsaut{6.0mm}
\VUfilet{22mm}
